%%%PAGE 016 rot=0 legibility=good summary=磁场中粒子圆轨迹技巧：平移圆、时间极值、单边有界入射与出射、汇聚发散、径入径出
%%%BLOCK id=016-01 ch=11 sec=平移圆 kind=method alt=none cont=none y=0.10-0.20
\textbf{平移圆}\quad $V_0$定\quad 入射位置改变
%FIG p016_a 0.10 0.12 0.36 0.195
\notefig[0.3]{figures/p016_a.png}
%%%END
%%%BLOCK id=016-02 ch=11 sec=时间极值 kind=method alt=none cont=none y=0.19-0.30
\textbf{时间极值}\quad $t=\dfrac{\theta m}{qB}=\dfrac{s}{V_0}$\quad 时间最长\ $\theta$最大\ 圆心角最大
% CHECK: 分子手写潦草，像 θm（也可能是 2πm）；"时间最长"后面一字同样像 θ（也可能是 弧）
\begin{itemize}
\item $V_0$定值\quad 最长弦/弧（劣）
\item $V_0$变化\quad 最大圆心角
\end{itemize}
% 原稿：以上两行左侧有大括号；"最长弦/弧（劣）"下方有一条波浪线向右延伸至"最大圆心角"上方
%%%END
%%%BLOCK id=016-03 ch=11 sec=单边有界入射 kind=method alt=none cont=none y=0.30-0.42
\textbf{单边有界入射}
%FIG p016_b 0.05 0.335 0.205 0.415
\notefig[0.25]{figures/p016_b.png}
角入角出\quad 速偏角$=$圆心角$=2$倍弦切角
%FIG p016_c 0.49 0.315 0.65 0.390
\notefig[0.25]{figures/p016_c.png}
\red{圆心圆}\quad \red{定速不定向问题}
% 原稿：红色弧线从"圆心圆"（浅粉底）右上角向下弯至"弦切角"右侧，连接这两处
%FIG p016_d 0.675 0.315 0.80 0.381
\notefig[0.25]{figures/p016_d.png}
垂直入射即半圆
%%%END
%%%BLOCK id=016-04 ch=11 sec=单边有界出射型 kind=method alt=none cont=none y=0.42-0.60
\textbf{单边有界出射型}\quad \red{O离板距$d$}\quad \red{光屏}
%FIG p016_e 0.22 0.415 0.52 0.52
\notefig[0.35]{figures/p016_e.png}
到不了\unclear{边界}交点，\red{到切点，已撞板了}
% CHECK: "到不了"与"交点"之间被光屏横线盖住的字看不清，像 边界/切
% 原稿：图右侧有一个竖直大括号，旁注 3m
$3\unit{m}$\quad $r=2.5\unit{m}$
%FIG p016_f 0.825 0.42 0.96 0.48
\notefig[0.25]{figures/p016_f.png}
% 红色小图（直角三角形）文字：斜边 R，竖直边 d-R（符号像 + 或 -，按物理意义取 -），水平边 sqrt(R^2-(d-R)^2)
% CHECK: 红字 d 与 R 之间的符号手写像 + 也像 -
\begin{itemize}
\item 最值相切：最远点为轨迹与边界切点
\item 最值相交：最远点为直径和边界交点
\item 能交则交，不能交则切。
\end{itemize}
% 原稿：以上三行左侧有大括号
%%%END
%%%BLOCK id=016-05 ch=11 sec=圆形磁场·汇聚发散 kind=method alt=none cont=none y=0.60-0.80
\textbf{汇聚\ 发散}
%FIG p016_g 0.02 0.64 0.52 0.815
\notefig[0.5]{figures/p016_g.png}
$V$大小变化\ 方向定

$B_1=B_2$\ 点入平出

平进点出
%%%END
%%%BLOCK id=016-06 ch=11 sec=径入径出 kind=method alt=none cont=none y=0.80-0.95
%FIG p016_h 0.04 0.795 0.27 0.97
\notefig[0.3]{figures/p016_h.png}
径入径出
%FIG p016_i 0.49 0.82 0.745 0.955
\notefig[0.3]{figures/p016_i.png}
化单边处理
%%%END
%%%PAGEEND 016
